\section{Related Work}
\label{sec:related}

\textbf{Deep phase picking.}
PhaseNet~\citep{zhu2019phasenet} predicts sample-wise P/S probabilities with a
U-Net-style architecture.
EQTransformer~\citep{mousavi2020eqt} jointly detects events and picks phases.
Generalized phase detection~\citep{ross2018gpd} and related CNN pickers
similarly map waveforms to phase scores.
Systematic benchmarks~\citep{munchmeyer2022benchmark,woollam2022seisbench}
show that absolute performance depends strongly on dataset shift and evaluation
protocol.

\textbf{Travel times and path corrections.}
Classical seismology uses distance-dependent travel-time tables and station or
path corrections estimated from residuals.
Our distance MLP plays the role of a smooth baseline; historical median
residuals supply a nonparametric path correction regularized by shrinkage.

\textbf{Catalog-assisted and historical picking.}
We deliberately restrict scope to settings with event context.
Methods that discover events from continuous waveforms (detection/association)
are complementary and out of scope.

\textbf{Candidate re-ranking.}
Peak finding on probability curves yields multiple S candidates; re-ranking
with external constraints is a lightweight alternative to fine-tuning the
picker.
Our ablations show that a fixed score already captures most of the gain
available from switching between PhaseNet and a history-informed pick among
existing candidates (oracle selector ceiling $\approx 0.720$ F1@$0.5$), while
oracle candidate selection ($0.890$) indicates remaining coverage limits.

\textbf{Uncertainty-aware picking.}
Some pickers expose probability curves or calibrated scores; we treat PhaseNet
probabilities as candidate features inside an explicit score rather than
claiming a new uncertainty model.
Additional uncertainty-calibration literature should be inserted only after
manual DOI verification (placeholder avoided here).
